% 编译方式: xelatex SL_h3_completeness_proof.tex
\documentclass[UTF8]{ctexart}
\usepackage{amsmath, amssymb, amsthm}
\usepackage[margin=2.5cm]{geometry}
\usepackage[hyphens]{url}
\usepackage[hidelinks]{hyperref}
\usepackage{booktabs}

\newtheorem{theorem}{定理}[section]
\newtheorem{proposition}[theorem]{命题}
\newtheorem{lemma}[theorem]{引理}
\newtheorem{corollary}[theorem]{推论}
\newtheorem{remark}[theorem]{注}
\newtheorem{definition}[theorem]{定义}

\title{第三左定空间 \texorpdfstring{$H^3[-1,1]$}{H3} 中多项式基的解析完备性: 证明文档}
\author{研究证明文档 (项目: BVE research)}
\date{2026-08-04}

\begin{document}
\maketitle

\begin{abstract}
本文证明移位 Krein Laplacian $K_c f = -f'' + c f$ ($c > 0$, 边界条件
$f'(\pm1) = \bigl(f(1)-f(-1)\bigr)/2$) 的第三左定空间
$(H^3[-1,1], (\cdot,\cdot)_{3,c})$ 中多项式基
\[
	p_0 = 1, \quad p_1 = x, \quad
	p_{2n} = x^{2n} - \tfrac{n}{n-1} x^{2n-2}, \quad
	p_{2n+1} = x^{2n+1} - \tfrac{n}{n-1} x^{2n-1} \quad (n \in \mathbb{N},\ n \geq 2)
\]
是解析完备的. 这是 Jones--Littlejohn--Quintero Roba [4, Section 6] 开放问题
在 $H^2$ 之后 (会话 9) 的自然推广. 证明初等且自足, 核心是新的 \emph{H 矩跳跃}
观察: 把 $H^1$-矩 $M_k := (w, x^k)_1$ 而非 $L^2$-矩作为基本量, 则正交条件
$(w, K_c p_n)_1 = 0$ 直接给出与 $H^2$ 情形相同的\emph{二阶}递推
$c M_{2m} = A_m M_{2m-2} - B_m M_{2m-4}$, 边界项 $w(1)\pm w(-1)$ 被完全吸收;
再利用已有的增长引理 ($M_{2m}$ 超阶乘增长) 与新的多项式增长上界
$|M_{2m}| \leq C\sqrt{m}$ (由 $w \in H^1$ 的 Cauchy-Schwarz 估计给出) 相矛盾,
迫使全部 $H^1$-矩为零, 从而 $w = 0$. 同一论证对一切整数 $s \geq 1$ 的左定空间
$H^s$ 成立 (推论 6.2). 文档末尾附对抗性审计清单与精确有理数数值验证.
\end{abstract}

\tableofcontents

\section{记号与问题}

\subsection{移位 Krein Laplacian 与左定空间族}

固定 $c > 0$. 在 $L^2(-1,1)$ 中定义
\begin{equation}
	K_c f := -f'' + c f, \qquad
	D(K_c) := \Bigl\{ f : f, f' \in AC[-1,1],\ f'' \in L^2(-1,1),\
		f'(\pm 1) = \tfrac{1}{2}\bigl( f(1) - f(-1) \bigr) \Bigr\}.
\end{equation}
由左定理论 (Littlejohn--Wellman [5,6]), 对 $s \geq 0$ 定义左定空间
\begin{equation}
	H^s[-1,1] := D\bigl(K_c^{s/2}\bigr), \qquad
	(f, g)_s := \bigl( K_c^{s/2} f,\ K_c^{s/2} g \bigr)_{L^2}.
\end{equation}
特别地 $H^0 = L^2$, $H^1 = D(K_c^{1/2})$, $H^2 = D(K_c)$, $H^3 = D(K_c^{3/2})$.
第三左定空间的内积可改写为
\begin{equation}
	(f, g)_3 = (K_c f, K_c g)_1 \qquad (f, g \in H^3),
\end{equation}
因为 $K_c^{1/2} K_c = K_c^{3/2}$ (泛函演算).

\subsection{问题陈述}

会话 9 已解决 $H^2$ 中的解析完备性. 本文推进到第三左定空间:

\begin{theorem}[主定理]
对每个 $c > 0$, 多项式序列 $\{p_n\}_{n \in \mathbb{N}_0, n \neq 2,3}$
在 Hilbert 空间 $(H^3[-1,1], (\cdot,\cdot)_{3,c})$ 中解析完备:
\[
	\overline{\operatorname{span}\{ p_n \}}^{H^3} = H^3[-1,1].
\]
\end{theorem}

证明由五步组成: (i) 等距同构 $K_c: H^3 \to H^1$ 把问题归结为 $\{K_c p_n\}$
在 $H^1$ 中的完备性; (ii) $H^1$-矩跳跃: 正交条件化为 $M_k = (w, x^k)_1$ 的
二阶递推; (iii) 增长引理 (会话 9 已证, 系数相同); (iv) 多项式增长上界
$|M_{2m}| \leq C\sqrt{m}$; (v) 矛盾迫使 $w = 0$.

\section{预备: 等距同构与 \texorpdfstring{$K_c p_n$}{Kc pn} 的显式形式}

\subsection{等距同构}

\begin{lemma}[谱与可逆性, 会话 9]
$K_c$ 自伴且具有紧预解式, $\sigma(K_c) = \{c\} \cup \{(n\pi)^2 + c\}
\cup \{\mu_n^2 + c\}$ ($\mu_n$ 为 $\tan\mu = \mu$ 的正根). 特别地 $0 \notin \sigma(K_c)$.
\end{lemma}

\begin{lemma}[等距同构]
$K_c: \bigl( H^3, (\cdot,\cdot)_3 \bigr) \to \bigl( H^1, (\cdot,\cdot)_1 \bigr)$
是等距线性同构, 且
\[
	\overline{\operatorname{span}\{ p_n \}}^{H^3} = H^3
	\iff
	\overline{\operatorname{span}\{ K_c p_n \}}^{H^1} = H^1 .
\]
\end{lemma}

\begin{proof}
等距性由 (3): $\|K_c f\|_1^2 = (K_c f, K_c f)_1 = (f, f)_3 = \|f\|_3^2$.
满射: 对 $g \in H^1 = D(K_c^{1/2})$, $K_c^{-1} g \in D(K_c^{3/2}) = H^3$
(泛函演算), 故 $K_c$ 的像覆盖 $H^1$; 结合等距性即单射. 等距同构保持
闭包与正交补, 故完备性等价. \qed
\end{proof}

\subsection{\texorpdfstring{$K_c p_n$}{Kc pn} 的显式形式}

\begin{lemma}[会话 9, 引理 4]
对 $n \geq 2$,
\begin{align}
	K_c p_{2n} &= c\, x^{2n} - A_n\, x^{2n-2} + B_n\, x^{2n-4},
	& A_n &= 2n(2n-1) + \frac{cn}{n-1}, \quad B_n = 2n(2n-3); \\
	K_c p_{2n+1} &= c\, x^{2n+1} - A'_n\, x^{2n-1} + B'_n\, x^{2n-3},
	& A'_n &= 2n(2n+1) + \frac{cn}{n-1}, \quad B'_n = 2n(2n-1).
\end{align}
\end{lemma}

\section{H 矩跳跃: 正交条件化为二阶递推}

本文的关键简化是: 对 $H^1$ 完备性问题, 取 $H^1$-矩而非 $L^2$-矩.

\begin{lemma}[H 矩递推]
设 $w \in H^1$ 且 $(w, K_c p_n)_1 = 0$ 对所有可允许的 $n$
($n \in \{0,1\} \cup \{4,5,\dots\}$) 成立. 记 $H^1$-矩
\begin{equation}
	M_k := (w, x^k)_1 = \int_{-1}^{1} w'(x^k)'\, dx + c \int_{-1}^{1} w x^k\, dx
		- \tfrac12 \Delta w\, \Delta(x^k),
	\qquad \Delta f := f(1) - f(-1).
\end{equation}
则 $M_0 = M_1 = 0$, 且对 $m \geq 2$ 有\emph{二阶}跳变递推
\begin{align}
	c\, M_{2m} &= A_m\, M_{2m-2} - B_m\, M_{2m-4}, \\
	c\, M_{2m+1} &= A'_m\, M_{2m-1} - B'_m\, M_{2m-3},
\end{align}
其中 $A_m, B_m, A'_m, B'_m$ 同引理 3.
\end{lemma}

\begin{proof}
$K_c p_0 = K_c 1 = c$, $K_c p_1 = c x$, 故
$(w, K_c p_0)_1 = c M_0 = 0$ 与 $(w, K_c p_1)_1 = c M_1 = 0$, 即 $M_0 = M_1 = 0$.
对 $m \geq 2$, 由引理 3 与内积对第二变元的线性:
\[
	0 = (w, K_c p_{2m})_1 = c (w, x^{2m})_1 - A_m (w, x^{2m-2})_1
		+ B_m (w, x^{2m-4})_1 = c M_{2m} - A_m M_{2m-2} + B_m M_{2m-4},
\]
即 (8). 奇次情形同理给出 (9). \qed
\end{proof}

\begin{remark}[边界项消失]
$L^2$-矩 $\mu_k = \int w x^k$ 满足三阶递推, 并携带边界项
$D = w(1)+w(-1)$ 与 $S = w(1)-w(-1)$ (会话 10 前段探索). 改用 $H^1$-矩后,
内积定义中的 $- \tfrac12 \Delta w\, \Delta(x^k)$ 项与分部积分边界项精确抵消,
递推降回二阶. 这是``矩跳跃''方法适用于任意左定空间的一般机制 (推论 6.2).
\end{remark}

\section{增长引理与多项式上界}

\subsection{增长引理 (会话 9)}

\begin{lemma}[增长引理]
设 $c > 0$, $u_0 = 0$, $u_1 = 1$, 且对 $j \geq 2$
\begin{equation}
	c\, u_j = A_j\, u_{j-1} - B_j\, u_{j-2},
	\qquad A_j = 2j(2j-1) + \frac{cj}{j-1}, \quad B_j = 2j(2j-3),
\end{equation}
则 $u_j > 0$, $u_j$ 单调不减, 且 $u_j \geq (4/c)^{j-1} j!$ 对一切 $j \geq 1$.
对系数 $(A'_j, B'_j)$ 同样成立.
\end{lemma}

\begin{proof}
见会话 9 证明文档 (引理 5): 单调性由 $A_j - B_j = 4j + \tfrac{cj}{j-1} \geq c$
归纳; 增长由比值 $r_j = u_j/u_{j-1} \geq 4j/c$ 连乘. \qed
\end{proof}

\subsection{多项式增长上界}

\begin{lemma}[H 矩的多项式界]
设 $w \in H^1$, $M_k = (w, x^k)_1$. 则存在常数 $C = C(c, \|w\|_{H^1})$ 使得对
一切 $m \geq 1$,
\begin{equation}
	|M_{2m}| \leq C \sqrt{m}, \qquad |M_{2m+1}| \leq C \sqrt{m}.
\end{equation}
\end{lemma}

\begin{proof}
由 (7) 展开. 偶次: $\Delta(x^{2m}) = 0$, 故
\[
	M_{2m} = 2m \int_{-1}^{1} w'(x) x^{2m-1}\, dx + c\, \mu_{2m},
	\qquad \mu_{2m} = \int_{-1}^{1} w(x) x^{2m}\, dx.
\]
由 Cauchy-Schwarz, $\|x^k\|_{L^2} = \sqrt{2/(2k+1)}$:
\[
	|M_{2m}| \leq 2m\, \|w'\|_2 \sqrt{\tfrac{2}{4m-1}}
		+ c\, \|w\|_2 \sqrt{\tfrac{2}{4m+1}}.
\]
因 $4m - 1 \geq 2m$ ($m \geq 1$), $2m / \sqrt{4m-1} \leq \sqrt{2m}$, 故
$|M_{2m}| \leq \bigl( \sqrt{2}\, \|w'\|_2 + c\|w\|_2 \bigr) \sqrt{2m} \leq C\sqrt{m}$.
奇次: $\Delta(x^{2m+1}) = 2$, $\Delta w = S$ 有界,
\[
	|M_{2m+1}| \leq (2m+1)\, \|w'\|_2 \sqrt{\tfrac{2}{4m+1}}
		+ c\, \|w\|_2 \sqrt{\tfrac{2}{4m+3}} + |S| \leq C\sqrt{m},
\]
因为 $(2m+1)^2 / (4m+1) \leq 2m+1$ 且 $|S| = |\int w'| \leq \sqrt2\, \|w'\|_2$.
\qed
\end{proof}

\section{主定理的证明}

\begin{proof}[主定理]
设 $f \in H^3$ 满足 $(f, p_n)_3 = 0$ 对所有可允许的 $n$ 成立. 由 (3) 与引理 2,
$w := K_c f \in H^1$ 满足
\[
	(w, K_c p_n)_1 = (K_c f, K_c p_n)_1 = (f, p_n)_3 = 0.
\]
令 $M_k = (w, x^k)_1$. 由引理 4, $M_0 = M_1 = 0$ 且 (8)(9) 成立.
偶矩序列由 $M_0 = 0$ 与自由参数 $M_2$ 唯一决定: $M_{2m} = M_2\, u_m$
($u_m$ 如引理 5). 若 $M_2 \neq 0$, 则
\[
	|M_{2m}| = |M_2|\, u_m \geq |M_2| \left( \frac{4}{c} \right)^{m-1} m! \to \infty
	\qquad (m \to \infty),
\]
与引理 6 的 $|M_{2m}| \leq C\sqrt{m}$ 矛盾. 故 $M_2 = 0$, 递推 (8) 给出全部
偶 $H^1$-矩为零. 奇次情形完全相同 (系数 $A'_m, B'_m$, 自由参数 $M_3$):
$M_3 = 0$, 全部奇 $H^1$-矩为零.

于是 $M_k = 0$ 对一切 $k \geq 0$ 成立, 即
\[
	(w, p)_1 = 0 \quad \text{对所有多项式 } p \text{ 成立}.
\]
多项式在 $H^1(-1,1)$ 中稠密 (Weierstrass 于 $C^\infty$ 稠密性 + $H^1$ 稠密性),
故取 $p \to w$ 得 $\|w\|_1^2 = (w, w)_1 = 0$, 即 $w = 0$. 由引理 2 的等距同构,
$f = K_c^{-1} w = 0$. 由 Hahn--Banach 推论, $\operatorname{span}\{p_n\}$
在 $H^3$ 中稠密. \qed
\end{proof}

\section{推广: 一切左定空间中的完备性}

\begin{corollary}[任意阶左定空间]
对每个整数 $s \geq 1$, 多项式基 $\{p_n\}$ 在 $(H^s, (\cdot,\cdot)_s)$ 中解析完备.
\end{corollary}

\begin{proof}[证明思路; 更正见 \texttt{SL\_denseness\_criteria.pdf} 定理 8]
原思路中恒等式
\[
	K_c^{s/2}p_{2m} = cK_c^{s/2}x^{2m} - A_mK_c^{s/2}x^{2m-2}
		+ B_mK_c^{s/2}x^{2m-4}
\]
对 $s \geq 4$ 不成立 (会话 11 数值证实: $K_c^2p_{2m}$ 有四项),
故下列证明是修正版.
第一步: 对任意整数 $t \geq 0$, $\{K_c p_n\}$ 在 $H^t$ 中完备. 设
$w \in H^t$ 且 $(w, K_c p_n)_t = 0$ 对一切允许 $n$ 成立, 记 $H^t$-矩
$N_k = (w, x^k)_t$. 由 $K_c p_{2m} = cx^{2m} - A_mx^{2m-2} + B_mx^{2m-4}$
(多项式恒等式, 与 $t$ 无关), 跳变递推 $cN_{2m} = A_mN_{2m-2} - B_mN_{2m-4}$
成立; 由 $\|x^k\|_t \leq C_t k^t$ 与 Cauchy-Schwarz, $|N_k| \leq Ck^t$
多项式增长, 与增长引理 5 的超阶乘增长矛盾, 故 $N_2 = N_3 = 0$, 全部
$H^t$-矩为零, $w = 0$.
第二步 (等距传输): $(f,g)_s = (K_c f, K_c g)_{s-2}$ 且 $K_c: H^s \to H^{s-2}$
等距同构 ($s \geq 2$), 故 $\{p_n\}$ 在 $H^s$ 中完备当且仅当 $\{K_c p_n\}$
在 $H^{s-2}$ 中完备. 迭代得 $\{p_n\}$ 在 $H^s$ 中完备当且仅当
$\{K_c^r p_n\}$ 在 $H^{s-2r}$ 中完备 ($0 \leq 2r \leq s$).
第三步 (基始): $H^0 = L^2$ 与 $H^1$ 由一阶矩准则
($\|x^k\|_{H^s} \leq Ck^{s-1/2}$, $s = 0, 1$ 时指数 $< 1$); $H^2, H^3$
由第一步 ($t = 0, 1$, 即会话 9 与本文主定理). 综合三步, 一切整数
$s \geq 0$ 成立. \qed
\end{proof}





\begin{remark}
$s = 1$ 即本文主定理 ($H^1$-矩路线); $s = 2$ 即会话 9 的 $H^2$ 结论
($L^2$-矩路线); 二者同构于同一矩跳跃机制. 对非整数 $s \geq 0$, 由
$H^s$ 的插值理论 (Fischbacher--Gesztesy--Hagelstein--Littlejohn [7]) 与
$\|x^k\|_s \leq C_s k^s$ 的相同估计, 结论对一切 $s \geq 1$ 亦成立
(实数情形本文不再展开).
\end{remark}

\section{数值验证}

脚本 \path{scripts/h3_v69b_h1moments.py} (精确有理数与浮点混合) 与
\path{scripts/h3_v68_bases_min.py} 验证以下内容.

\begin{enumerate}
	\item \textbf{H 矩递推恒等式 (精确)}: 对测试多项式 $w \in \{x^2, x^3,
		x^2+x^4, 1+x+x^5\}$ 与 $c = 3$,
		$c M_{2m} - A_m M_{2m-2} + B_m M_{2m-4} = (w, K_c p_{2m})_1$ 对
		$m = 2, \dots, 10$ 逐项精确相等 (偶/奇均验证).
	\item \textbf{增长引理 (精确)}: 对 $c \in \{1,3,10,50\}$, 序列
		$u_m$ ($u_1 = 1$) 满足 $u_m \geq (4/c)^{m-1} m!$ 到 $m = 30$;
		$c=3$ 时 $u_{20}$ 为 39 位整数.
	\item \textbf{矛盾量级}: $c=3$ 时 $(4/c)^{m-1}m! / (C\sqrt m)$ 在
		$m = 20$ 时约 $8\times10^{19}$; 矛盾远非边界情形.
	\item \textbf{$H^1$ 投影残差 (浮点)}: $x^2, x^3$ 到
		$\operatorname{span}\{K_c p_n : \deg \leq 2N-2\}$ 的 $H^1$ 投影残差平方
		($c=3$): $x^2$ 在 $\deg \leq 18$ 时 $1.1\times10^{-12}$, $\deg \leq 26$
		时低于机器精度; $x^3$ 在 $\deg \leq 18$ 时 $4.0\times10^{-14}$.
		这与 $\{K_c p_n\}$ 在 $H^1$ 中完备一致.
	\item \textbf{前段探索 (三阶系统, 未用于证明)}: 显式积分解
		$\prod(1 \pm 1/(2k))$ (偶), $\prod(1 + 3/(2k)), \prod(1+1/(2k))$ (奇)
		符号级验证为齐次解; 比值固定点 $e_j = 1+1/(2j)$ (偶),
		$1+3/(2j)$ (奇) 为精确不动点; 偏差 $d_j = \rho_j - e_j \geq 0$ 且
		单调递减 ($j \leq 2\times10^6$); 最小解 $h^*$ 满足 $h^*_0 \neq 0$.
		这些结果记录于研究总结文档, 因 H 矩路线而不再需要.
\end{enumerate}

\section{对抗性审计}

\begin{description}
	\item[逻辑闭环] 无循环论证: 等距同构 (引理 2) 使用谱结论 $0 \notin
		\sigma(K_c)$ (独立计算); 增长引理 (引理 5) 为会话 9 的独立引理;
		上界 (引理 6) 只使用 $w \in H^1$ 与 Cauchy-Schwarz; 完备性结论不
		被任何一步预支.
	\item[可允许指标] 递推只使用 $n = 0, 1$ 与 $n \geq 4$; 缺省次数 2, 3
		不参与论证 (与 $H^2$ 情形一致), 因此代数不完备性不影响解析完备性.
	\item[边界情形] $m = 1$ 时 $M_2, M_3$ 为自由参数, 恰为矛盾论证的起点;
		$c > 0$ 是 $0 \notin \sigma$ 的本质条件; $B_m \geq 0$ 与
		$A_m - B_m \geq c$ 对全部 $m \geq 2$ 已验证.
	\item[H 矩良定性] $M_k = (w, x^k)_1$ 对 $w \in H^1$ 良定 ($w' \in L^2$,
		$w \in L^2$, $\Delta w$ 存在), $x^k$ 光滑.
	\item[多项式稠密性] 多项式在 $H^1(-1,1)$ 中稠密为标准结论 (先以
		$C^\infty$ 逼近 $w$, 再以多项式一致逼近导数).
	\item[结论强度] 上界 (11) 是多项式增长; 增长引理是超阶乘 ($(4/c)^{m-1}m!$).
		两者差距极大, 结论不依赖精细常数.
\end{description}

\section{涉及到的数学知识}

\begin{description}
	\item[左定理论] 对正自伴算子 $A$, $H_s = D(A^{s/2})$ 连同
		$(f,g)_s = (A^{s/2}f, A^{s/2}g)$ 构成左定空间族; $H^3 = D(K_c^{3/2})$
		即 $s = 3$ 的成员 [5,6,7].
	\item[等距同构传递] $K_c: H^3 \to H^1$ 保内积双射, 把稠密性问题在
		源/目标空间之间传递; $M^\perp = \{0\}$ 等价于 $\overline M = H$
		(Hahn-Banach 推论).
	\item[H 矩跳跃] 对与正交条件同源的内积取矩 (本文: $H^1$-矩), 边界项
		被内积中的边界项吸收, 递推阶数降低 (三阶 $\to$ 二阶). 这是
		``矩跳跃'' 方法的关键机制.
	\item[二阶递推支配解] $c u_j = A_j u_{j-1} - B_j u_{j-2}$ 的支配解行为由
		系数差 $A_j - B_j$ 控制; 单调性 + 比值下界给出阶乘增长.
	\item[Cauchy-Schwarz 与矩界] $L^2$ 函数 $g$ 的矩满足
		$|\langle g, x^k \rangle| \leq \|g\|_2 \sqrt{2/(2k+1)}$; 本文对
		$H^1$-矩用 $\int w' x^{k-1}$ 项给出 $\sqrt m$ 界.
	\item[Weierstrass 与 Sobolev 稠密性] 多项式在 $C$、$L^2$、$H^1$ 中稠密;
		矩全零 $\Rightarrow$ 函数为零.
	\item[泛函演算] $K_c^{1/2}K_c = K_c^{3/2}$ 在 $D(K_c^{3/2})$ 上成立,
		支撑 (3) 与等距同构.
\end{description}

\begin{thebibliography}{9}
\bibitem{lq} L. L. Littlejohn, A. Quintero, \emph{Krein-Sobolev orthogonal
	polynomials}, in: Special Functions and Orthogonal Polynomials (OPSFA-16),
	Springer, CRM Series in Mathematical Physics (2025), 107--123.
	\url{https://link.springer.com/chapter/10.1007/978-3-031-90135-5_7}
\bibitem{axioms} A. Jones, L. L. Littlejohn, A. Quintero Roba, \emph{Krein-Sobolev
	Orthogonal Polynomials II}, Axioms 14 (2025), no. 2, 115.
	\url{https://doi.org/10.3390/axioms14020115}
\bibitem{lw1} L. L. Littlejohn, R. Wellman, \emph{A general left-definite theory
	for certain self-adjoint operators with applications to differential equations},
	J. Differential Equations 181 (2002), 280--339.
	\url{https://doi.org/10.1006/jdeq.2001.4079}
\bibitem{lw2} L. L. Littlejohn, R. Wellman, \emph{On the spectra of left-definite
	operators}, Complex Anal. Oper. Theory 7 (2013), 437--455.
	\url{https://doi.org/10.1007/s11785-011-0214-1}
\bibitem{fg} C. Fischbacher, F. Gesztesy, P. Hagelstein, L. L. Littlejohn,
	\emph{Abstract left-definite theory: A model operator approach, examples,
	fractional Sobolev spaces, and interpolation theory}, arXiv:2408.01514 (2024).
	\url{https://arxiv.org/abs/2408.01514}
\bibitem{h2proof} 项目文档, \emph{第二左定空间 $H^2[-1,1]$ 中多项式基的解析完备性:
	证明文档}, 2026-08-04 (会话 9), \texttt{docs/SL\_h2\_completeness\_proof.pdf}.
\bibitem{weier} K. Weierstrass, \emph{Approximation theorem} (经典 Weierstrass 定理).
\end{thebibliography}

\end{document}

